% Drawing 87/02-3 -- folder 87, batch 02, archive page 27 (PDF page 28).
% Transcription: en marge gauche, deux croquis. En haut, un segment oblique épais portant trois points marqués, avec les lettres $x$ (près du premier) et $y$ (près du troisième). En bas, un triangle aplati : à gauche les points $x'$ et $x$, en haut $y''$, $y$ et $y'$, à droite $z$ et $z'$ ; des traits joignent $x$ à $y$ et à $z$, $y$ à $z$, $y'$ à $z'$, et un trait fin $x'$ à $y''$.
% Kind: segment with marked points; flattened triangle x, y, z with outer points (transitivity). Method 2 (schematic redraw from the facsimile). Status: ok.
% Notes: Top: the heavy oblique segment, three marked points on it (computed on the segment), the letter x under the first and y under the third. Bottom: the points x', x (left), y'', y, y' (top), z, z' (right) with their letters, and the strokes x'-x (heavy), x-y continued to y', x-z, y-z, y'-z', the thin stroke x'-y'' and a short stroke from y'' towards y. Left out: a small sign beside z (a struck letter or a second z), a short stroke under y (the tail of its letter). Record checked against its .tex header after the rebuild of 7b758b8 (all fields agree). Overlaid 2026-11-12 (Opus 5.5) on the 400 dpi colour crop and corrected: the lower figure lay 2.5 to 3.5 units too high at the left and right (x', x, z, z'); its seven points are moved onto the blots of the page and every stroke redrawn between them; on the page the upper stroke starts at x and the right-hand strokes end at z (from y) and z' (from y'), which settles the two readings left open before. The transcription gives this figure as a row of labelled points in a displayed formula and describes it in the note; the mark is added beside it. Redrawn 2026-11-12 (issue 128): to check again.
% Coordinates: 1 unit = 1 mm = 1/100 of the crop width (bbox in page fractions: [0.16, 0.4, 0.42, 0.59]).
% Redrawn by Claude (claude-fable-5-2), issue #128, 2026-11-12; not checked by a human; not compiled (no TeX in the session).
% Overlaid by Claude (claude-opus-5-5) 2026-11-12 (issue #128): rasterised from its own coordinates over the 400 dpi colour facsimile and corrected from it. Not compiled. Status: ok.
% Rights: reproduces part of the work -- behind RENDER_TRANSCRIPTS, never published by itself.
\documentclass[tikz,border=3pt]{standalone}
\usepackage{amsmath,amssymb}
\usetikzlibrary{patterns,arrows.meta}
\tikzset{pencil/.style={black!45}, struck/.style={black!35}, mpt/.style={circle,draw,inner sep=1.3pt}, pt/.style={circle,fill,inner sep=0.9pt}, opt/.style={circle,draw,fill=white,inner sep=1.1pt}}
\begin{document}
\begin{tikzpicture}[x=1mm,y=1mm,line cap=round,line join=round,font=\scriptsize]
\draw[very thick] (9.5,-35.5) -- (45,-12);
\node[pt] at (16,-31.2) {};
\node[pt] at (29,-22.6) {};
\node[pt] at (36,-18) {};
\node[inner sep=1pt] at (18.5,-31.5) {\small $x$};
\node[inner sep=1pt] at (39.5,-22) {\small $y$};
\draw[very thick] (13.5,-84.3) -- (18.8,-81);
\draw (18.8,-81) -- (40.5,-72.5) -- (49,-69.3);
\draw (18.8,-81) -- (67,-88.5);
\draw (40.5,-72.5) -- (67,-88.5);
\draw (49,-69.3) -- (75.5,-92);
\draw[pencil] (13.5,-84.3) -- (35,-71);
\draw (35,-71) -- (38.5,-72.3);
\node[pt] at (13.5,-84.3) {};
\node[pt] at (18.8,-81) {};
\node[pt] at (35,-71) {};
\node[pt] at (40.5,-72.5) {};
\node[pt] at (49,-69.3) {};
\node[pt] at (67,-88.5) {};
\node[pt] at (75.5,-92) {};
\node[inner sep=1pt] at (9,-87.5) {\small $x'$};
\node[inner sep=1pt] at (21,-84.5) {\small $x$};
\node[inner sep=1pt] at (32,-67.5) {\small $y''$};
\node[inner sep=1pt] at (40.5,-78) {\small $y$};
\node[inner sep=1pt] at (54,-68.5) {\small $y'$};
\node[inner sep=1pt] at (64.5,-92) {\small $z$};
\node[inner sep=1pt] at (79.5,-92) {\small $z'$};
\end{tikzpicture}
\end{document}
